
\section{Open Problems}
\label{openProblems}
Graph adversarial learning has many problems worth to be studied by observing existing researches. In this section, we try to introduce several major research issues and discuss the potential solutions. Generally, we will discuss the open problems in three parts: attack, defense and evaluation metric.

\subsection{Attack}
We mainly approach the problems from three perspectives, attacks' side effects, performance, and prerequisites. First, we hope the inevitable disturbances in the attack can be stealthy and unnoticeable. Second, now that data is huge, attacks need to be efficient and scalable. Finally, the premise of the attack should not be too ideal, that is, it should be practical and feasible. Based on the limitations of previous works detailed in Section \ref{attack_limitation} and the discussion above, we rise several open problems in the following:

\begin{itemize}
	\item \textbf{Unnoticeable Attack}.  Adversaries want to keep the attack unnoticeable to avoid censorship. To this end, Z\"ugner et al. \cite{zugner2019adversarial} argue that the main property of the graph should be marginally changed after attack, e.g., attackers are restricted to maintain the degree distribution while attacking a specified graph, which is evaluated via a test static. Nevertheless, most of existing works ignore such a constraint despite achieving outstanding performance. To be more concealed, we believe that attackers should focus more on their perturbation impacts on a target graph, with more constraints placed on the attacks.
	\item \textbf{Efficient and Effective Algorithms}. Seeing from Table \ref{attack_table}, the most frequently used benchmark datasets are rather small-scale graphs, e.g., Cora \cite{mccallum2000automating} and Citeseer \cite{sen2008collective}. Currently, the majority of proposed methods are failed to attack a large-scale graph, for the reason that they need to store the entire information of the graph to compute the surrogate gradients or loss even with small changes, leading to the expensive computation and memory consumption. To address this problem, Z\"ugner et al. \cite{zugner2019adversarial} make an effort that derives an incremental update for all candidate set (potential perturbation edges) with a linear variant of GCN, which avoid the redundant computation while remain efficient attack performance. Unfortunately, the proposed method isn't suitable for a larger-scale graph yet. With the shortcoming that unable to conduct attacks on a larger-scale graph, a more efficient and effective attack method should be studied to address this practical problem. For instance, given a target node, the message propagation are only exits within its $k$-hops neighborhood.\footnote{Here $k$ depends on the specified method of message propagation and is usually of a smaller value, e.g., $k=2$ with a two-layer GCN.} In other words, we believe that adversaries can explore a heuristic method to attack an entire graph (large) via its represented subgraph (small), and naturally the complexity of time and space  are reduced.
	\item \textbf{Constraints on Attackers' Knowledge}. To conduct a real and practical black-box attack, attackers should not only be blind to the model knowledge, but also be strict to minimal data knowledge. This is a challenging setting from the perspective of attackers. Chen et al. \cite{chen2017practical} explore this setting by randomly sampling different size of subgraphs, showing that the attacks are failed with a small network size, and the performance increases as it gets larger. However, few works are aware of the constraints on attackers' knowledge, instead, they assume that perfect knowledge of the input data are available, and naturally perform well in attacking a fully ``exposed'' graph. Therefore, several works could be explored with such a strict constraint, e.g., attackers can train a substitute model on a certain part of the graph, learn the general patterns of conducting attacks on graph data, thus transfer to other models and entire graph with less prior knowledge.
	\item \textbf{Real-world Attack}. Attacks can't always be on an ideal assumption, i.e., studying attacks on a simple and static graphs. Such ideal, undistorted input is unrealistic in real cases. In other words, how can we improve existing models so that they can work in complex real-world environments?
\end{itemize}

\subsection{Defense}
Besides the open problems of attack mentioned above, there are still many interesting problems in defense on the graph, which deserve more attentions. Here we list some open problems that worth discussing:
\begin{itemize}

	\item \textbf{Defense on Various Tasks}. From table \ref{defense_table} we can see that, most existing works of defense \cite{zugner2019certifiable,xu2019topology,jin2019power,wu2019adversarial,wang2019adversarial} are focusing on the node classification task on graph, achieving promising performance. Except for node classification, there are various tasks on graph domain are important and should be pay more attention to. However, only a few works \cite{pezeshkpour2019investigating,zhou2019adversarial} try to investigate and improve model robustness on link prediction task on graph while a few works \cite{hou2019alphacyber,ioannidis2019graphsac,he2018adversarial} make an effort to defense on some tasks (e.g., malware detection, anomaly detection, recommendation) on other graph domains. Therefore, it is valuable to think about how to improve model robustness on various tasks or transfer the current defense methods to other tasks.

	\item \textbf{Efficient Defense.} Training cost is critical important factor to be taken into consideration under the industrial scenes. Therefore, it's worth to be studied how to guarantee acceptable cost while improving robustness. However, the currently proposed defense methods have rarely discussed the space-time efficiency of their algorithms. Wang et al. \cite{wang2019adversarial} design a bottleneck mapping to reduce the dimension of input for a better efficiency, but their experiments absence the consideration of large-scale graph. Z\"ugner et al. \cite{zugner2019certifiable} test the number of coverage iterations of their certification methods under the different number of nodes, but the datasets they used are still small. Wang et al. \cite{wang2019graphdefense} do not restrict the regularized adjacency matrix to be discrete during the training process, which improves the training efficiency.

	\item \textbf{Certification on Robustness}. Robustness of the graph model is always the main issue among all the existing works including attack and defense mentioned above. However, most researchers only focus on improving model (e.g., GCN) robustness, and try to prove the model robustness via the performance results of their model. The experimental results can indeed prove the robustness of the model to a certain extent. However, considering that the model performance is easily influenced by the hyperparameters, implementation method, random initialization, and even some hardware constraints (e.g., cpu, gpu, memory), it is difficult to guarantee that the promising and stable performance can be reobtained on different scenarios (e.g., datasets, tasks, parameter settings). There is a novel and cheaper way to prove robustness via certification that should be taken more seriously into consideration, which certificate the node's absolute robustness under arbitrary perturbations, but currently only a few works have paid attention to the certification of GNNs' robustness \cite{zugner2019certifiable,bojchevski2019certifiable} which are also a valuable research direction.

\end{itemize}

\subsection{Evaluation Metric}
As seen in Section \ref{metrics}, we got so many evaluation metrics in graph adversarial learning field, however, the number of effectiveness-relevant metrics are far more than the other two, which reflects the research emphasis on model performance is unbalanced. Therefore, we propose several potential works can be further studied:

\begin{itemize}
	\item \textbf{Measurement of Cost.} There are not many metrics to explore the model's efficiency which reflects the lack of research attention on it. The known methods roughly measure the cost via the number of modified edges which begs the question that is there another perspective to quantify the cost of attack and defense more precisely? In real world, the cost of adding an edge is rather different from removing one,\footnote{Usually, removing an edge is more expensive than adding one.} thus the cost between modified edges are unbalanced, which needs more reasonable evaluation metrics.

	\item \textbf{Measurement of Imperceptibility.} Different from the image data, where the perturbations are bounded in $\ell_p$-norm \cite{sun2018adversarial,goodfellow2014explaining,madry2017towards}, and could be used as an evaluation metric on attack impacts, but it is ill-defined on graph data. Therefore, how to better measure the effect of perturbations and make the attack more unnoticeable?
	\item \textbf{Appropriate Metric Selection.} With so many metrics proposed to evaluate the performance of the attack and defense algorithms, here comes a problem that how to determine the evaluation metrics in different scenarios?
\end{itemize}


\section{Conclusion}
\label{conclusion}
In this survey, we conduct a comprehensive review on graph adversarial learning, including attacks, defenses and corresponding evaluation metrics. To the best of our knowledge, this is the first work systemically summarize the extensive works in this field. Specifically, we present the recent developments of this area and introduce the arms race between attackers and defenders. Besides, we provide a reasonable taxonomy for both of them, and further give a unified problem formulation which makes it clear and understandable. Moreover, under the scenario of graph adversarial learning, we summarize and discuss the major contributions and limitations of current attack and defense methods respectively, along with the open problems of this area that worth exploring. Our works cover most of the relevant evaluation metrics in the graph adversarial learning field as well, aiming to provide a better understanding on these methods. In the future, we will measure the performance of proposed models with relevant metrics based on extensive empirical studies.

Hopefully, our works will serve as a reference and give researchers a comprehensive and systematical understanding of the fundamental issues, thus become a well starting point to study in this field.
